\documentclass[11pt]{article}
\usepackage[margin=2.6cm]{geometry}
\usepackage{amsmath,amssymb}
\usepackage{bm}
\usepackage[T1]{fontenc}
\usepackage[utf8]{inputenc}
\usepackage{parskip}
\usepackage{microtype}
\newcommand{\neff}{n_{\mathrm{eff}}}
\newcommand{\R}[1]{\medskip\noindent\textbf{#1}\quad}
\newenvironment{rquote}{\begin{quote}\itshape}{\end{quote}}
\begin{document}
\begin{center}
{\LARGE\bfseries Response to Reviewers}\\[6pt]
{\large Manuscript MATG-D-26-00180 --- \emph{Mathematical Geosciences},
second revision}\\[4pt]
``Effective-sample-size calibration of a modality test for cluster homogeneity
and the number of classes in spatially autocorrelated data''\\[4pt]
Francisco Parra \,\textperiodcentered\, Universidad de Santiago de Chile
\end{center}
\medskip

I thank Reviewer~2 for a second reading and for closing the file, and
Reviewer~1 for a report that is, once again, more useful than comfortable.
Reviewer comments are quoted in italics; section, equation and table numbers
below refer to the revised manuscript.

Three structural changes answer the report's opening paragraph, and much of
the detail follows from them.

\emph{A supplement.} As asked, twice, there is now a supplementary document. It
carries the full discussion of what Proposition~1 does and does not establish,
together with its proof (Section~S1); a complete specification of all sixteen
benchmarked criteria and their tuning constants (Section~S2); and an
implementation description that replaces, and considerably exceeds, the
appendix the reviewer found uninformative (Section~S3); and the comparison of
the fitted integral range and the simulation-based estimator against the
default, which was Section~4.5 of the reviewed version (Section~S4). The
appendix is gone from the paper, and Section~4.5 is now a summary paragraph.

\emph{A relativised framing.} The paper no longer suggests, anywhere I could
find it, that a clustering is uniquely determined or that the standard criteria
are wrong on spatial data. This was the report's most valuable criticism and it
has changed the abstract, the introduction, Section~2, Section~3.2, the
paragraph headings in Section~4.3 and the conclusions.

\emph{Less promised.} Claims of guarantee, verification, exactness and reliable
negatives are removed throughout, and Section~3.9 now opens by stating that no
result covers the deployed procedure.

The paper is the same length as the version reviewed, at thirty-eight pages,
despite the explanations added in response to the detailed points below; the
appendix, the proof, the discussion of the Proposition's scope and the study of
the simulation-based estimator have moved into the eight-page supplement.

\section*{General points}

\R{G1. There is no uniquely true clustering.}
\begin{rquote}
``The paper implicitly suggests in many places that there is a uniquely true
clustering of a dataset. This is not the case, see Hennig (2015). \dots\ any
wording that suggests that clusterings can only be valid if spatial
autocorrelation is taken into account should be changed (starting from the
abstract `miscalibrated on spatial data'). Indexes like the Silhouette index
may well be fine for spatial data if the application doesn't require what is
required here.''
\end{rquote}
I accept this without reservation, and it has improved the paper. The claim
that the standard criteria are miscalibrated on spatial data was doing work it
could not support, and removing it makes the contribution easier to state, not
harder. The abstract now says that those criteria ``take no account of the
dependence between observations, so on spatial data they leave one question
unanswered: whether an apparent grouping reflects discrete classes, or is
already produced by the autocorrelation of a single class.''

The introduction states the position explicitly, in a new paragraph:
\begin{rquote}
``None of this makes those criteria wrong. A clustering is not a fact waiting
to be uncovered: which groupings count as clusters depends on what the analysis
is for (Hennig 2015), and an analyst may have good reason to partition an
autocorrelated field into contiguous units whether or not the partition could
be explained by the autocorrelation alone. What these criteria cannot do is
answer one particular question, and it is the question this paper addresses:
does a region contain discrete classes, over and above the variation that a
single autocorrelated class already produces?''
\end{rquote}
Hennig (2015) is now cited there. The specific instances the reviewer names are
gone: ``a criterion that cannot recognise the absence of clusters is a
precarious foundation'' has been deleted; the conclusions no longer speak of
``inventing classes on featureless fields'' but of reporting several classes on
fields simulated with one; and the paragraph heading ``Standard indices
over-detect without exception'' is now ``Standard indices report classes on
every class-free scene.'' Where the simulation results are described, the claim
is now about what was simulated rather than about what exists: the criteria
report more than one class on scenes generated without discrete classes. That
is a statement about a known generating model, which is the only version of the
claim I am entitled to make.

\R{G2. Two-means may cut through a cluster when more than two are present.}
\begin{rquote}
``I don't think the authors have commented on the possibility that 2-means may
cut through a cluster in situations with more than 2 clusters in an early stage
of the algorithm.''
\end{rquote}
Correct, and the omission mattered, because this failure is not symmetric with
the other one. Section~3.3 now says so:
\begin{rquote}
``If an early split cuts through a class, the two fragments of that class are
handed to different subtrees, and no later test can reunite them: the damage is
permanent, and shows up as an over-estimate of $K$ or as a partition that
misplaces boundaries even when $\hat K$ is right. What the recursion does
repair is the opposite error, a split that leaves two classes on the same side,
since the child is tested again.''
\end{rquote}
The passage goes on to note that this is one reason the paper reports partition
quality alongside $\hat K$, and that recovery holding for $K$ up to five bounds
how often the failure is severe at these separations without showing that it
does not occur. The corresponding limitation in Section~6 is stated in the same
terms.

\section*{Detailed points}

\R{1. Fixing $K$ before any classifier runs.}
\begin{rquote}
``An analyst mapping land cover from a satellite scene must fix K before any
unsupervised classifier runs'' -- not true in general. There are clustering
methods (such as DBSCAN) that estimate K on the fly.''
\end{rquote}
Corrected. The sentence now reads that the analyst ``must commit to $K$ for the
large family of methods that take it as input, among them $k$-means and finite
mixtures; methods that instead determine the number of groups on the fly, such
as DBSCAN (Ester et al.\ 1996), replace that choice by tuning parameters that
govern it.''

\R{2. Swapping in a better null.}
\begin{rquote}
``Correcting this is not a matter of swapping in a better null.'' Isn't this
what you are doing (at least in some sense)? I don't really see what's wrong
with ``swapping in a better null''.
\end{rquote}
The sentence was wrong as written, since the paper does swap the null. What I
meant, and now say, is that swapping the null is not \emph{sufficient} while
the statistic stays the same, and the reason is worth spelling out: if the
reference is matched to the covariance estimated from the data, the
between-class contrast is itself part of that covariance, so the reference
reproduces the observed dispersion whether or not classes are present. The text
now reads ``the null must change, but so must the statistic'', with the
argument given in full in Section~2.3 (see also point~9).

\R{3. Terminology introduced and then varied.}
\begin{rquote}
``Terms like `scene', `mosaic', `MiniBatch' are used without explicit
introduction \dots\ Later you use terms like `$\ell^2$ proxy', `$1/e$ proxy'
again without having explicitly defined what you mean.''
\end{rquote}
Section~3.1 now defines both terms at first use: ``we call one such image over
$D$ a \emph{scene}, and the spatial pattern the classes form across it the
\emph{mosaic}''. The two names for the same object have been reduced to one:
the manuscript now says ``the $\ell^{2}$ proxy'' everywhere, and ``$1/e$''
appears only where the crossing itself is described. MiniBatch $k$-means is
introduced with a citation in Section~4.2 and specified in Section~S2.

\R{4. Reliable negatives.}
\begin{rquote}
``its negatives are trustworthy: when it declares a region homogeneous that
verdict can be relied on''. Not so. In general if a statistical test doesn't
reject a null hypothesis, it never means that the null hypothesis is true
\dots\ don't write as if it were different.
\end{rquote}
Removed, both in the introduction and in the abstract, and I am glad to be rid
of it. The introduction now reads: ``Because the test is tuned for specificity,
it seldom divides a region that the null can account for; as with any test,
failing to reject is not evidence that a region is homogeneous, only that the
data do not contradict homogeneity at the effective sample size.'' The abstract
speaks of the trade-off rather than of trustworthy verdicts, and the
conclusions of ``the asymmetry between its two kinds of error''.

\R{5, 29, 30, 32. What the Proposition does not establish.}
\begin{rquote}
``As Proposition 1 is `stylised', it doesn't guarantee anything for the
proposed method. There is no such thing as a proven `asymptotic
false-split-control property' for the method as a whole.'' \dots\ ``The
result is deliberately stylised'' -- this is sugarcoating the fact that you
apparently don't have a better one \dots\ The result is of purely heuristic
value.
\end{rquote}
The reviewer is right, and I have stopped hedging around it. Section~3.9 now
opens: ``We have no result covering the deployed procedure. What follows
concerns a stylised version of the test \dots\ It is a consistency check on the
calibration principle rather than a statement about the method as applied.''
The phrase ``asymptotic false-split-control property'' no longer appears in the
abstract, the introduction, or the conclusions; each says instead that for a
stylised version of the test we prove that the probability of dividing a
homogeneous region vanishes. Section~3.9 also states, where the reviewer asked
for it, that ``nothing above establishes a property of the procedure as
deployed; its adequacy is a matter of evidence''. The sentence claiming the
calibration ``is designed so that \dots\ the recursion does not split'' has
been replaced by one saying what is intended and what is actually shown.

On whether to keep the Proposition at all, which the reviewer left to me: I
have kept it, but moved its proof and the detailed discussion of its scope to
Section~S1 and cut Remark~1 from fifty-seven lines to sixteen. My reason for
keeping it is that the reduction it performs is the only place where the
paper's central premise is stated in a form that can be checked, and a reader
who doubts (C3) can now see precisely where the doubt bites. My reason for
moving it out is the reviewer's, that it does not deserve the space it took.

\R{6. Which criteria assume independence.}
\begin{rquote}
``None of these criteria was proposed with a distribution theory that requires
independent observations.'' Well, some yes (mixture/BIC, gap), some no (ASW,
DB). Be precise.
\end{rquote}
That sentence was simply wrong and has been replaced by a precise one in
Section~2.1: the gap statistic and model-based selection by BIC do rest on
assumptions treating observations as independent, the gap through its reference
distribution and the BIC through a likelihood that is a product over
observations; the silhouette, Calinski--Harabasz and Davies--Bouldin indices
assume no probability model at all, ``and on dependent data they are neither
correct nor incorrect, because they make no probabilistic claim that dependence
could violate.'' The same distinction is now made in Section~3.2 and in
Section~6, where the phrase ``the independence reference of the classical
criteria'' has been corrected (point~47).

\R{7. Spatial mixture models.}
\begin{rquote}
``By the way, there is also mixture model-based clustering taking spatial
autocorrelation into account, see, e.g., Vanhatalo et al. (2021).''
\end{rquote}
Added in Section~2.1, immediately after the mixture and BIC sentence.

\R{8. A reference is a model, not a method.}
\begin{rquote}
``That an independence reference over-detects structure in dependent data is
known across several fields.'' Be precise. Structure is detected by methods for
detecting structure. A reference is not a method but a model.
\end{rquote}
Reworded: ``That methods comparing data against an independence reference
report structure which the dependence alone can account for is known across
several fields.''

\R{9. The dispersion argument.}
\begin{rquote}
``Because k-means dispersion is itself a variance-based quantity, a null
matched to the data's covariance structure can explain away the very clusters
one seeks'' -- I don't understand the logic.
\end{rquote}
It was underexplained. Section~2.3 now gives the argument in steps: keep a
dispersion-based index and improve only the null, matching it to the covariance
estimated from the data; ``the contrast between classes contributes to the very
covariance the null now reproduces; the null therefore generates reference
datasets of much the same dispersion as the data whether or not classes are
present, and the comparison loses what it was meant to detect.'' Hence a
criterion able to survive such a null must read something the covariance does
not determine.

\R{10. Not all the listed criteria use a reference distribution.}
\begin{rquote}
``Internal validity criteria for K, namely the gap statistic \dots\ the average
silhouette, Calinski-Harabasz and Davies-Bouldin, compare a within-cluster
dispersion against the value expected with no clustering, under an independent
reference distribution.'' This is true for the gap statistic but not for the
other criteria listed here.
\end{rquote}
Corrected in Section~3.2, which now attributes the reference distribution to
the gap alone and describes the other three as optimising a ratio of within- to
between-cluster scatter over $K$ with no reference distribution at all.

\R{11. There are no independent observations to reduce.}
\begin{rquote}
``autocorrelation reduces the number of independent observations far below N''
-- as a standard there aren't any independent observations here.
\end{rquote}
Reworded: ``autocorrelation means the sample carries the information of far
fewer than $N$ independent observations''.

\R{12. Drift and dependence are separate failures.}
\begin{rquote}
``You could have an independence null that takes drift into account but not
dependence.''
\end{rquote}
Section~3.2 now says exactly that: the two failures ``are distinct, and a
reference could model the drift while still ignoring the dependence, or the
reverse.''

\R{13, 15, 18, 35. Small definitional points.}
``The null $H_0$'' is now ``the null hypothesis $H_0$''. ``Bimodality'' has
been replaced by ``a genuine departure from unimodality'' where more than two
modes are possible. The tile index $j$ is introduced explicitly in Section~3.5.
The separation $\delta$ is defined where the power panel first uses it, as the
standard deviation of the class mean vectors introduced in Section~4.1.

\R{14. Where Assumption 1 is needed.}
\begin{rquote}
``Chances are you should state already here that you sometimes require
Assumption 1 and sometimes not.''
\end{rquote}
A paragraph now follows the assumption saying which parts of the paper use it:
the default ESS-Dip uses only unimodality, through the distribution-free
uniform calibration, whereas the Gaussian form licenses ESS-Dip-R, the
indicator integral range, the simulation reference and condition (C1). Each of
those places repeats the dependence.

\R{16. In what sense the uniform is least favourable.}
\begin{rquote}
``least-favourable $>$ least favourable (also elsewhere) By the way \dots\ It
is unclear here in which sense the uniform is least favourable.''
\end{rquote}
The sense is now stated where the term is introduced: the uniform ``is least
favourable among unimodal laws in the sense that it maximises the asymptotic
distribution of the rescaled dip, so that a critical value taken from it bounds
the level for the others''. The predicative uses are unhyphenated. I should add
that the previous revision already narrowed this claim: Hartigan and Hartigan
prove it for laws with exponentially decreasing densities and conjecture it in
general, and the paper says so.

\R{17. The two ends of one calibration.}
\begin{rquote}
``the two are the distribution-free and the model-matched ends of one
calibration.'' Unclear.
\end{rquote}
Replaced by a plain statement of what each protects against: ``The two differ
only in the reference law, and so in what they protect against: ESS-Dip against
any within-class law in the class for which the uniform bound is proved,
ESS-Dip-R against the Gaussian of Assumption~1 specifically.''

\R{19, 21. The 27 per cent.}
\begin{rquote}
``27\%? Wasn't it supposed to be the majority?'' \dots\ ``There are only two
numbers `above', one of which is 27\% (for the non-straddling tiles \dots)!?''
\end{rquote}
The reviewer caught a real inconsistency: the text asserted that most tiles fall
within a single class and then reported a figure showing that on one of the
three scenes they do not. Section~3.5 now says so plainly, and explains why the
estimator survives it:
\begin{rquote}
``We check it, and it holds on some scenes but not all: 62\% of tiles are
single-class on Salinas and on the Sentinel-2 image, against 27\% on the
fragmented Indian Pines fields. Where the premise fails the estimate does not
collapse, because a tile straddling a boundary is not thereby uninformative:
most of its area still lies in one class, and its $1/e$ crossing is lengthened
rather than replaced by the scene scale. On Indian Pines the tile median still
returns 3.5 pixels against a scene-wide estimate of 16. We therefore treat the
premise as a condition under which the estimate is clean, not one the method
requires.''
\end{rquote}
I have removed the appeal to the median's breakdown point that stood here
before, since with 27 per cent single-class tiles the contaminated fraction
exceeds one half and the argument does not apply.

\R{20. ``Exact for the exponential model''.}
\begin{rquote}
``I don't understand the use of `exact' here -- a definition that is not exact
is not a proper definition.''
\end{rquote}
The word was attached to the wrong noun. The passage now separates the
convention from the model: ``Taking the $1/e$ crossing as the range is a
convention: it reproduces the range parameter of the exponential model, for
which $\rho(h)=\exp(-h/\ell)$ gives $\rho(\ell)=e^{-1}$, and approximates it for
the other models.''

\R{22, 23. Equation (8).}
\begin{rquote}
``(8): How is $r_j^2=\int r\rho(r)dr$, the latter apparently not dependent on
j?'' \dots\ ``The fitted parameters'' -- fitted how?
\end{rquote}
The correlogram in the integral is the model fitted to tile $j$ and does depend
on $j$; it is now written $\rho_j$ and identified as such in the following
line. The fitting is now specified in place: the parameters are fitted ``by
ordinary least squares to the tile's empirical variogram with the model of
smallest residual sum of squares retained''.

\R{24. Simulating from a Gaussian is not assumption-free.}
\begin{rquote}
``Simulating from a Gaussian is for sure not `assumption-free'.''
\end{rquote}
Agreed and corrected: the alternative is now described as one ``that dispenses
with the geometric constant, at the price of assuming a Gaussian model
explicitly''.

\R{25. A smooth drift is not unimodal.}
\begin{rquote}
``A smooth drift $\mu$ is itself unimodal'' -- how so?
\end{rquote}
That was a category error. What matters is that the drift adds no modes while
still distorting the range estimate, and the text now says that: ``A smooth
drift contributes no modes of its own: it displaces the feature vectors
continuously across the scene rather than gathering them at distinct locations.
Left in place, however, it lengthens the estimated range $\ell$, deflating
$\neff$ and biasing the test against splitting.''

\R{26. ``The trend world''.}
\begin{rquote}
``from the trend world of the factorial design'' -- I don't understand this.
\end{rquote}
The term was used before its definition. It now reads ``on the class-free
scenes carrying a drift that Section~4.1 calls the \emph{trend} world''.

\R{27. Two meanings of $B$.}
\begin{rquote}
``B in Sec. 3.7 and 3.8 is not the same thing.''
\end{rquote}
Correct. The simulations per node in Section~3.7 are now $M$; $B$ is reserved
for the ensemble size, as in the notation table.

\R{28. The $k$-means ensemble.}
\begin{rquote}
``k-means is about optimising the k-means objective function, so if you get B
solutions from B initialisations, some of these are good and some aren't. You
shouldn't just use the median or mean of K derived from these. The way I think
it should be done is that you run B initialisations for each K and then for
each K you use the best solution according to the objective function.''
\end{rquote}
This is done, and the fault is mine for not saying so: within each node the
binary split already retains the best of five random starts by the $k$-means
objective, exactly as the reviewer describes. The ensemble is a second and
different thing, and Section~3.8 now distinguishes them:
\begin{rquote}
``The choice among $k$-means initialisations \emph{within} a node is made in
the usual way, by the $k$-means objective: each binary split retains the best
of five random starts. What the objective cannot settle is which of several
near-equivalent splits the recursion then follows, since a different split at
one node changes the nodes its descendants see. To absorb that residual path
dependence the recursion is repeated $B$ times.''
\end{rquote}
Section~S3 repeats the distinction where the constants are listed.

\R{31. The consequence does not follow from the bound.}
\begin{rquote}
``Consequently \dots'' This does not follow from (13) unless $\alpha$ is 0.
\end{rquote}
Quite so; it follows from the sharper statement in the same sentence, not from
the bound. The Proposition now says: ``and, under the same conditions, that
limit is in fact zero. It is this last statement, not the bound (13), which
yields the following \dots''

\R{33. Verification.}
\begin{rquote}
``Empirical verification'' is not a verification in any generality and
shouldn't be called ``verification''. I'd accept ``validated''.
\end{rquote}
All three uses are changed: the tile premise is now ``checked'', the empirical
level is ``validated'', and the simulation study ``measures how far the
calibration keeps the recursion from splitting''. The one remaining use of
``verification'' refers to verifying condition (C3) mathematically, which is
the sense the reviewer would allow.

\R{34. Simulation-based tests do not control levels exactly.}
\begin{rquote}
``A simulation-based test does not control levels `exactly' unless you have
infinitely many simulation runs.''
\end{rquote}
Corrected in both places: the reference now controls the level ``up to Monte
Carlo error, and up to the fitted variogram''.

\R{36. Figure 1(b) depends on the alternative.}
\begin{rquote}
``this will depend on the exact alternative, which at the very least needs
specification of $\delta$ if not more.''
\end{rquote}
The separation is now given both in the text and in the caption, with the
dependence stated: ``the curve is specific to that separation and moves with
it.''

\R{37. References in Section 3.10.}
\begin{rquote}
``Give references for the claims made here.''
\end{rquote}
Added: Cooley and Tukey (1965) for the transform, Sculley (2010) for the
mini-batch update, and Hartigan and Hartigan (1985) for the cost of the dip.

\R{38, 39, 41. Wording.}
``the spatial gap next'' is now ``the autocorrelation-matched gap described
next''; ``the respective later work'' is now ``the work that introduced each of
those calibrations, not \dots\ the papers that introduced the indices
themselves''; and the heading now reads ``Comparison with the method that
ESS-Dip extends.''

\R{40. Scaled to which standard deviation?}
\begin{rquote}
``scaled to each band's standard deviation'' -- true or estimated, if
estimated how? If true, does this make sense given that in reality this
information isn't available?
\end{rquote}
Estimated, and the text now says so and why: ``as estimated from the scene
itself, since no deployable version of the criterion has access to the true
value''. Section~S2 gives the construction, including the bandwidth
$\sigma=\ell/2$ at which the generator's correlogram crosses $1/e$ at $\ell$.

\R{42. Tuning choices for all methods.}
\begin{rquote}
``In a Supplement there needs be an explanation of all tuning choices for
*all* methods simulated \dots\ (``scenes subsampled to 1200 pixels'' -- how
exactly? Reference? ``Bonferroni correction'' was how exactly applied? Note
that as a standard this needs a prespecification of the number of tests run,
and I don't think you can do that as it is not clear in advance when the
splitting process will end.)
\end{rquote}
Section~S2 of the supplement now specifies every criterion: the partitions and
their initialisations, the gap's reference construction and one-standard-error
rule, the subsample sizes and why the indices requiring pairwise distances need
them (the subsample is uniform without replacement, drawn once per scene rather
than per $K$, and shared across the three indices), the calibrated validity
index with its $B=20$ null datasets and threshold, the two nulls that
distinguish its variants, and the ablations.

On the Bonferroni correction the reviewer's objection is well taken and the
supplement concedes it before answering it. Prespecification is possible here
only because the recursion is capped: a binary recursion returning at most
$K_{\max}$ leaves performs at most $K_{\max}-1$ splits and therefore visits at
most $2K_{\max}-1$ nodes, so $\alpha/K_{\max}$ is a valid, if conservative,
level for the family of tests actually run. The supplement states plainly that
without the cap the number of tests would not be bounded in advance and the
correction would have no prespecified denominator.

\R{43. Table 4.}
\begin{rquote}
``Say explicitly what the entries are.''
\end{rquote}
The caption now identifies each row: the proportion of one-class scenes divided
at all, the mean of $\hat K$ over the four-class scenes, and the proportion of
those on which $\hat K$ equalled four.

\R{44. Must the autocorrelation be real?}
\begin{rquote}
``The calibration presumes a within-class autocorrelation that is real'' -- is
that so? In principle everything should go through if the autocorrelation is
zero, in which case $\neff=n$, shouldn't it?
\end{rquote}
It should, and the sentence was misleading. Section~5.2 now opens: ``Nothing in
the calibration requires the within-class autocorrelation to be non-zero: were
it absent, $\ell$ would be one pixel, $a=1$ and $\neff=n_C$, and the test would
reduce to the ordinary dip test at the nominal size. What the calibration does
require is that the range be estimable and separable from the between-class
mosaic, since it is the \emph{ratio} $n_C/a$ that the test consumes.''

\R{45. ``The true number of classes is one''.}
\begin{rquote}
``if the `true number of classes' is defined according to (1), there is no
guarantee that in a real situation with a single land-cover class the true
number of classes is indeed one \dots\ You can say something like `in such a
situation the desirable outcome is one class'.''
\end{rquote}
Adopted, with the reasoning made explicit: ``the desirable outcome is a single
class. The phrase is deliberate: the ground truth records one land-cover label,
which need not mean that the model of Equation~(1) holds there with $K=1$, and
a window whose dominant class covers less than all of its pixels could
legitimately support a second. What we can say is that a criterion reporting
several classes on such a window is reporting structure the label map does not
recognise.''

\R{46. References for ARI and NMI.}
\begin{rquote}
``Give a reference for ARI and NMI. (I think ARI is enough by the way.)''
\end{rquote}
Both suggestions adopted. Hubert and Arabie (1985) is now cited for the
adjusted Rand index, and normalised mutual information has been dropped from
Table~8 and from the surrounding discussion, which loses nothing: the two
indices ranked the methods identically on the hyperspectral scenes, and the
conclusion drawn from the table, that a test which refuses to split does not
maximise partition recovery, rests on the ARI alone.

\R{47. The independence reference of the classical criteria.}
Corrected in Section~6, which now says that Assumption~1 ``replaces the
independence reference of the gap statistic, and supplies one where the ratio
indices have none''.

\R{48, 49. The conclusions.}
The closing paragraph has been rewritten and shortened, and no longer restates
the effective-sample-size argument for a third time. ``Inventing classes on
featureless fields'' is gone, along with the rest of the absolute language, as
described under G1.

\bigskip
\noindent
I am grateful for the care in this report. Several of its points were things I
had written without noticing that I could not support them, and the paper is
more honest for their removal.

\end{document}
